% Repaired diplomatic French. Decisions: TOH-P119-R1-DIP-001 through -004.
% Supersedes the active use of TOH-DIP-P119-001 without overwriting it.
\SourcePage{119}
\LayerTarget{tohoku:paper}
\begin{center}
  {\Large\bfseries SUR QUELQUES POINTS D'ALGÈBRE HOMOLOGIQUE.\par}
  \vspace{1.2em}
  {\large\scshape Alexandre Grothendieck\par}
  \vspace{1.2em}
  {(Reçu mars 1, 1957)\par}
\end{center}

\LayerTarget{tohoku:introduction}
\LayerContents{Introduction}
\section*{\centering Introduction\textsuperscript{1)}}

\LayerTarget{tohoku:introduction:I}
\LayerTarget{tohoku:introduction:I:p1}
\paragraph{I. Contenu du travail.}
Ce travail a son origine dans une tentative d'exploiter l'analogie formelle
entre la théorie de la cohomologie d'un espace à coéfficients dans un faisceau
[4], [5] et la théorie des foncteurs dérivés de foncteurs de modules [6], pour
trouver un cadre commun permettant d'englober ces théories et d'autres.

\LayerTarget{tohoku:introduction:I:p2}
Ce cadre est esquissé dans le Chapitre I, dont le thème est le même que celui
de [3]. Ces deux exposés cependant ne se recouvrent pas, sauf dans le seul
N\textsuperscript{o} 1.4. Je me suis attaché notamment à donner des critères
maniables, à l'aide de la notion de sommes et produits infinis dans les
catégories abéliennes, pour l'existence de ``suffisamment'' d'objets injectifs
ou projectifs dans une catégorie abélienne, sans quoi les techniques
homologiques essentielles ne peuvent s'appliquer. De plus, pour la commodité
du lecteur, une place assez large a été faite à l'exposé du langage fonctoriel
(N\textsuperscript{os} 1.1, 1.2 et 1.3). L'introduction des catégories
additives au N\textsuperscript{o} 1.3, préliminaire aux catégories abéliennes,
fournit un langage commode (par exemple pour traiter des foncteurs spectraux
au Chapitre II).

\LayerTarget{tohoku:introduction:I:p3}
Le Chapitre II esquisse les points essentiels du formalisme homologique dans
les catégories abéliennes. La parution de [6] m'a permis d'être très concis,
les techniques de Cartan--Eilenberg se transportant sans aucun changement dans
le nouveau cadre. Les numéros 2.1 et 2.2 ont été écrits cependant de façon à ne
pas exclure les catégories abéliennes ne contenant pas assez d'objets injectifs
ou projectifs. Dans les numéros suivants, nous employons à fond les techniques
usuelles de résolutions. Les N\textsuperscript{os} 2.4 et 2.5 contiennent des
compléments divers et sont essentiels pour la compréhension de la suite. En
particulier, le théorème 2.4.1 donne une façon mécanique d'obtenir la plupart
des suites spectrales connues (et en tous cas \emph{toutes} celles rencontrées
dans ce travail).

\LayerTarget{tohoku:introduction:I:p4}
Dans le Chapitre III nous redéveloppons la théorie de la cohomologie d'un
espace à coéfficients dans un faisceau, y inclus les suites spectrales
classiques de Leray. L'exposé donné ici représente un assouplissement par
rapport à [4], [15], en particulier en ce que tous les résultats essentiels
sont obtenus sans faire, à presque aucun moment dans ce Chapitre (pas plus que
dans les suivants), d'hypothèse restrictive sur la nature des espaces
envisagés; de sorte

\LayerTarget{tohoku:introduction:footnote:1}
\SourceFootnote{1}{L'essentiel des Chapitres I, II, IV et une partie du Chapitre III a été
développé au printemps 1955, à l'occasion d'un séminaire d'Algèbre Homologique
donné à l'Université de Kansas. ---Les numéros entre crochets renvoient à la
fin de cet article.}

\TohokuContinuation{120}

